\begin{styletable}

Равносильность по определению
    & $A \eqdef B$, $C(x) \eqdefr D(x)$
    & \verb|\eqdef|, \verb|\eqdefr|
\\

Отрицание
    & $\lnot A$
    & \verb|\lnot A|
\\

Конъюнкция \newline(И)
    & $A \land B$
    & \verb|A \land B|
\\

Дизъюнкция \newline(ИЛИ)
    & $A \lor B$
    & \verb|A \lor B|
\\

Следование (импликация)
    & $A \im B$
    & \verb|A \im B|
\\

Эквивалентность
    & $A \eq B$
    & \verb|A \eq B|
\\

Квантор общности
    & $\forall x \in X \qs x > 2$
    & \verb|\forall x \in X \qs x > 2| 
\\

Квантор существования
    & $\exists x \in X \qs x = 3$
    & \verb|\exists x \in X \qs x = 3|
\\

Единственность существования
    & $\existso x \in X \qs x = 3$
    & \verb|\existso x \in X \qs x = 3|
\\

\hline
\end{styletable}

\stylenote{
    Команды \texttt{\textbackslash eqdef} ($\eqdef$) и \texttt{\textbackslash eqdefr} ($\eqdefr$) используются исключительно для определения новых логических свойств и утверждений.
    Для определения числовых величин или множеств применяется равенство по определению: команды \texttt{\textbackslash is} ($\is$) и \texttt{\textbackslash isr} ($\isr$).
}

\stylenote{
    При использовании кванторов после указания переменной (и, возможно, ограничений на неё, отделяемых запятой) всегда ставится команда \texttt{\textbackslash qs}, которая выводит двоеточие, которое является чисто пунктуационным и логической нагрузки не несёт.
    Пример сложного выражения с кванторами:
    \begin{equation*}
        \forall x, y \in X, x \le y \qs
        \exists a, b \in \R \qs
            a + b > x - y.
    \end{equation*}
}